\chapter{Reprodução da análise e compilação}
\label{ap:reproducao}

A análise utiliza Python 3.13 e Matplotlib 3.10.7. Na raiz da implementação completa identificada no Apêndice~\ref{ap:rastreabilidade}, os comandos abaixo leem os dados arquivados e geram os artefatos, sem executar experimentos. Esses scripts e o pacote derivado estão preservados na branch \texttt{master}; não integram a branch de compilação do Overleaf:

\begin{lstlisting}[style=tcccode,xleftmargin=3.4em,xrightmargin=0.4em,framexleftmargin=0pt,framexrightmargin=0pt,framesep=3pt,numbersep=9pt]
python -B prova-conceito/benchmark/m01_analyze.py E11-M01-auditoria-extracao.zip
python -B prova-conceito/benchmark/m01_publication.py
python -B docs/tcc2/gerar_ilustracoes_monografia.py
\end{lstlisting}

O primeiro comando valida o pacote e gera métricas e gráficos; o segundo gera as tabelas de apresentação e os contrastes. O terceiro desenha os diagramas da arquitetura, que não são capturas de novas execuções. A versão revisada da figura de camadas é gerada por \path{uteis/gerar_camadas_ambiente.py}, na branch \texttt{tcc2/overleaf-utfpr}.

A Tabela~\ref{tab:arquivos-pacote} descreve o pacote derivado. As durações brutas estão em nanossegundos, convertidos conforme a Seção~\ref{sec:metricas}.

\begin{table}[htbp]
\centering\small
\caption{Conteúdo do pacote derivado usado na análise.}
\label{tab:arquivos-pacote}
\begin{tabularx}{\linewidth}{@{}p{4cm}X@{}}
\toprule
Arquivo & Conteúdo \\
\midrule
\texttt{inventario.json} & Identidades, revisões, hashes e vínculos com os arquivos de cada rodada. \\
\texttt{rodadas.csv} & 30 linhas com cenário, bloco, posição e duração monotônica do lote. \\
\texttt{requisicoes.csv} & 30.000 registros com rodada, passo, identidade, status e duração individual. \\
\texttt{checksums.sha256} & Hashes dos três arquivos para verificação de integridade. \\
\bottomrule
\end{tabularx}
\fonte{}
\end{table}

Para compilar a monografia, o arquivo principal é \texttt{main.tex}. O projeto utiliza pdfLaTeX, BibTeX, a classe \texttt{utfpr.cls} e o estilo \texttt{abntex2-alf.bst}. As figuras estão em \texttt{figuras/}; as tabelas derivadas, em \texttt{capitulos/dados-e11-m01/}; os trechos de código são incluídos diretamente dos arquivos identificados. Uma sequência de compilação local é:

\begin{lstlisting}[style=tcccode,xleftmargin=3.4em,xrightmargin=0.4em,framexleftmargin=0pt,framexrightmargin=0pt,framesep=3pt,numbersep=9pt]
pdflatex -interaction=nonstopmode -halt-on-error main.tex
bibtex main
pdflatex -interaction=nonstopmode -halt-on-error main.tex
pdflatex -interaction=nonstopmode -halt-on-error main.tex
\end{lstlisting}

No Overleaf, deve-se importar a estrutura do projeto sem reorganizar as pastas, selecionar \texttt{main.tex} como documento principal e pdfLaTeX como compilador. A folha de aprovação oficial deve ser incorporada aos pré-textuais após sua emissão.
